\documentclass{standalone}
% translate with >> pdflatex -shell-escape <file>

% This file is an extract of the PGFPLOTS manual, copyright by Christian Feuersaenger.
% 
% Feel free to use it as long as you cite the pgfplots manual properly.
%
% See
%   http://pgfplots.sourceforge.net/pgfplots.pdf
% for the complete manual.
%
% Any required input files (for <plot table> or <plot file> or the table package) can be downloaded
% at
% http://www.ctan.org/tex-archive/graphics/pgf/contrib/pgfplots/doc/latex/
% and
% http://www.ctan.org/tex-archive/graphics/pgf/contrib/pgfplots/doc/latex/plotdata/

\usepackage{pgfplots}
\usepackage{pgfplotstable}
\pgfplotsset{compat=newest}

% The manual loads these libraries globally; the extracted standalone examples
% need them too. external/clickable/snakes are intentionally omitted.
\usetikzlibrary{
	arrows.meta,backgrounds,calc,decorations.footprints,decorations.markings,
	decorations.pathreplacing,fixedpointarithmetic,intersections,patterns,
	plotmarks,shapes.arrows,spy,through}
\usepgfplotslibrary{
	colorbrewer,colormaps,dateplot,fillbetween,groupplots,patchplots,polar,
	smithchart,statistics,ternary,units}

% Some colormap examples reuse this display style defined earlier in the manual.
\pgfplotsset{
	cycle from colormap manual style/.style={
		x=3cm,y=10pt,ytick=\empty,
		colorbar style={x=,y=,ytick=\empty},
		point meta min=0,point meta max=1,
		stack plots=y,
		y dir=reverse,colorbar style={y dir=reverse},
		every axis plot/.style={line width=2pt},
		legend entries={0,...,20},
		legend pos=outer north east,
	},
}

\pagestyle{empty}

\begin{document}
\begin{tikzpicture}
\begin{axis}
    % FokkerDrI_layer_0.facetIdx.dat contains:
    % # each row makes up one facet; it
    % # consists of 0-based indices into
    % # the vertex array
    % 0    1    2 % triangle of vertices #0,#1 and #2
    % 0    3    1 % triangle of vertices #0,#3 and #1
    % 3    4    1
    % 5    6    7
    % 6    8    7
    % 8    9    7
    % 8    10   9
    % ...
    % while FokkerDrI_layer_0.vertices.dat contains
    % 105.577    -19.7332    2.85249    % vertex #0
    % 88.9233    -21.1254    13.0359    % vertex #1
    % 89.2104    -22.1547    1.46467    % vertex #2
    % 105.577    -17.2161    12.146
    % 105.577    -10.6054    18.7567
    % 105.577    7.98161     18.7567
    % 105.577    14.5923     12.146
    % ...
    \addplot3 [patch,shader=interp,
        patch table=
            {plotdata/FokkerDrI_layer_0.facetIdx.dat}]
        file
            {plotdata/FokkerDrI_layer_0.vertices.dat};
\end{axis}
\end{tikzpicture}
\end{document}
